\documentclass[11pt,a4paper]{article}

% ---------- Packages ----------
\usepackage[utf8]{inputenc}
\usepackage[T1]{fontenc}
\usepackage[margin=1in]{geometry}
\usepackage[protrusion=true,expansion=false]{microtype}
\usepackage{setspace}
\usepackage{amsmath, amssymb}
\usepackage{booktabs}
\usepackage{tabularx}
\usepackage{array}
\usepackage{enumitem}
\usepackage{titlesec}
\usepackage{xcolor}
\usepackage[hidelinks,colorlinks=true,linkcolor=blue!50!black,citecolor=blue!50!black,urlcolor=blue!50!black]{hyperref}
\usepackage{fancyhdr}
\usepackage{parskip}

% ---------- Spacing ----------
\onehalfspacing

% ---------- Section formatting ----------
\titleformat{\section}{\Large\bfseries\color{blue!40!black}}{\thesection}{0.6em}{}
\titleformat{\subsection}{\large\bfseries\color{blue!30!black}}{\thesubsection}{0.6em}{}
\titlespacing*{\section}{0pt}{1.4ex plus 1ex minus .2ex}{0.8ex plus .2ex}
\titlespacing*{\subsection}{0pt}{1.0ex plus .8ex minus .2ex}{0.6ex plus .2ex}

% ---------- Header / Footer ----------
\pagestyle{fancy}
\fancyhf{}
\fancyhead[L]{\small\itshape Literature Review \& Research Gap}
\fancyhead[R]{\small\thepage}
\renewcommand{\headrulewidth}{0.3pt}

% ---------- Custom commands ----------
\newcolumntype{Y}{>{\raggedright\arraybackslash}X}
\newcommand{\paperref}[1]{\textit{#1}}

% ---------- Title ----------
\title{
  \vspace{-1.5cm}
  \rule{\linewidth}{0.4pt}\\[0.4em]
  {\Huge\bfseries Literature Review and Research Gap}\\[0.3em]
  {\Large Inventory Management under Carbon Emission Regulations \\ and Time-Varying Demand}\\[0.4em]
  \rule{\linewidth}{0.4pt}
}
\author{}
\date{\today}

\begin{document}
\maketitle
\thispagestyle{empty}

% ===== Abstract =====
\begin{abstract}
\noindent
This document synthesises the findings of five seminal works that lie at the
intersection of \emph{inventory control}, \emph{carbon-emission regulation}, and
\emph{time-varying demand}. The reviewed papers fall into two largely
disjoint streams: (i) carbon-aware inventory models, in which classical
Economic Order Quantity (EOQ) frameworks are extended to incorporate
emission caps, carbon taxes, cap-and-trade markets, and green-technology
investment; and (ii) inventory models with non-stationary
demand patterns, particularly the \emph{power demand pattern}. After
critically appraising each contribution, the document identifies the gaps
that remain at the confluence of the two streams and proposes specific
directions for future research.
\end{abstract}

\vspace{0.3em}

% ============================================================
\section{Introduction}
% ============================================================

The intensifying global focus on climate change has placed
\emph{operational} levers of carbon abatement at the centre of
sustainable supply-chain research. Regulators across the European Union,
North America, and Asia have implemented a variety of carbon-pricing
mechanisms, including carbon taxes, mandatory emission caps, and
cap-and-trade markets, while consumer pressure increasingly rewards firms
that publicise green-technology investments and sustainability
campaigns. In parallel, the classical inventory-control literature has
relaxed the constant-demand assumption of the standard EOQ model and
explored several time-varying demand structures, of which the
\emph{power demand pattern} captures both products whose demand
intensifies near the end of the cycle (e.g., petrol, staple
groceries) and products whose demand wanes (e.g., perishables).

Each of these two streams has matured independently, yet their
intersection, namely \emph{carbon-constrained inventory control under
non-stationary, power-form demand}, remains comparatively unexplored.
This review summarises five representative works that anchor each
stream, exposes their methodological assumptions, and articulates the
research gap that motivates further investigation.

% ============================================================
\section{Overview of the Reviewed Papers}
% ============================================================

Table \ref{tab:overview} provides a high-level comparison of the five
reviewed studies along the dimensions of demand structure, regulatory
framework, decision variables, and central methodological contribution.

\begin{table}[h!]
\centering
\caption{Comparative overview of the five reviewed papers.}
\label{tab:overview}
\renewcommand{\arraystretch}{1.25}
\small
\begin{tabularx}{\textwidth}{@{}p{3.6cm}p{2.4cm}p{2.7cm}YY@{}}
\toprule
\textbf{Reference} & \textbf{Demand} & \textbf{Carbon Policy} & \textbf{Decision Variable(s)} & \textbf{Core Contribution} \\
\midrule
Hua, Cheng \& Wang (2011)\textsuperscript{$\dagger$}    & Constant       & Cap-and-trade                                  & Order quantity                                       & Closed-form EOQ extension under carbon trading; analyses effect of cap and price. \\
Benjaafar, Li \& Daskin (2013)                          & Constant       & Strict cap, tax, cap-and-offset, cap-and-trade & Lot size, multi-firm coordination                    & Lot-sizing models; quantifies value of operational adjustment vs.\ technology investment and supply-chain collaboration. \\
Hasan, Roy, Daryanto \& Wee (2021)                      & Linear in green tech.\ \& promotion & Carbon tax, cap-and-trade, strict cap & Order quantity, green-tech investment             & Three integrated EOQ models with green investment and demand-stimulating promotion. \\
Sicilia, Gonz\'alez-De-la-Rosa, Febles-Acosta \& Alcaide-L\'opez-de-Pablo (2014) & Power pattern  & None                                      & Reorder point, scheduling period, lot size           & Algorithmic optimal policy for production-inventory system with backlogged shortages and proportional production rate. \\
\bottomrule
\end{tabularx}
\\[0.3em]
{\footnotesize\textsuperscript{$\dagger$} The SSRN working paper (\texttt{ssrn1628953}) is the pre-print of Hua et al.\ (2011) and is therefore treated as a single contribution.}
\end{table}

% ============================================================
\section{Thematic Literature Review}
% ============================================================

\subsection{Stream 1: Carbon-aware Inventory Management}

\subsubsection*{Hua, Cheng \& Wang (2011)}

Hua, Cheng and Wang formulate one of the earliest closed-form EOQ
extensions in which a retailer operates under a cap-and-trade emission
regime. The authors explicitly model carbon emissions arising from two
operational sources: (i)~transportation of replenishment orders, taken to
be linear in the number of orders placed per unit time, and (ii)~warehouse
storage, taken to be linear in the average inventory level. By appending the
purchased or sold carbon credits to the classical EOQ cost function, they
derive an optimal order quantity $Q^{*}$ and prove a clean structural
result that compares $Q^{*}$ to the cost-minimising EOQ quantity $Q^{0}$
and the emission-minimising quantity $\hat{Q}$.

A central insight of the paper is the role played by the ratio
$g/e$ (warehouse-to-transport carbon intensity) versus the ratio
$h/K$ (warehouse-to-transport monetary intensity). When the two ratios
coincide, the cost-minimising order quantity simultaneously minimises
emissions, so participation in the carbon market is unnecessary; when
they differ, $Q^{*}$ falls strictly between $Q^{0}$ and $\hat{Q}$. The
analytical and numerical results highlight that carbon trading
\emph{mechanically links} cost and emissions, but they rest on the strong
assumption of \emph{deterministic, time-invariant} demand and an
exogenous carbon price.

\subsubsection*{Benjaafar, Li \& Daskin (2013)}

Benjaafar, Li and Daskin generalise the carbon-aware paradigm beyond a
single firm and beyond cap-and-trade, building on the multi-period
single- and multi-stage lot-sizing models that underlie much of the
operations-management literature. Their contribution is fourfold.
First, they show that emission parameters can be associated with each
decision variable in classical lot-sizing problems, yielding tractable
formulations under four distinct policies: strict caps, carbon tax,
cap-and-offset, and cap-and-trade. Second, by varying these parameters
numerically, they generate twelve managerial observations, of which the
most striking is that \emph{operational adjustment alone}, with no
investment in greener technology, can deliver substantial emission
reductions at modest cost. Third, they isolate the conditions under
which tighter caps may paradoxically increase aggregate emissions,
thereby informing policy design. Fourth, they evaluate the value of
\emph{supply-chain collaboration}, demonstrating that the gains depend
sensitively on the prevailing regulatory mechanism and on whether
emission caps are imposed firm-wise or supply-chain-wide.

The analytical core of the paper, an EOQ-style appendix that derives a
piecewise-defined optimal order quantity under a strict cap, complements
the multi-period programming results and unifies the single-firm
treatment of Hua et al.\ (2011) with multi-firm extensions. Like its
predecessor, however, the model presumes deterministic, stationary
demand.

\subsubsection*{Hasan, Roy, Daryanto \& Wee (2021)}

The work of Hasan and colleagues represents the most recent and
arguably the most behaviourally rich contribution to the
carbon-aware EOQ literature reviewed here. They formulate three
integrated profit-maximisation models that simultaneously optimise the
order quantity $Q$ and the green-technology investment $G$ under (i)~a
carbon tax, (ii)~cap-and-trade, and (iii)~a strict carbon limit. Two
features distinguish their formulation. First, the green investment is
modelled by the concave function $R(G)=aG-bG^{2}$, capturing
diminishing returns of abatement technology. Second, demand is modelled as
$D(G,v)=D_{0}+D_{1}R(G)+mv$, so that demand is stimulated both by
emission-reduction effort and by an explicit promotion level $v$.

Their numerical and sensitivity analyses produce three policy-relevant
findings: (i)~the carbon \emph{tax rate} significantly drives total
profit, whereas the cap level under cap-and-trade or under a strict
limit is relatively benign; (ii)~promotion is profitable across all
three regimes; and (iii)~green-technology investment is profitable
because it contracts the emissions exposure that must be priced through
trading or limited by the cap. The model nevertheless remains rooted in
a single buyer's EOQ structure with a \emph{linear} demand function,
i.e.\ no explicit time variation within a cycle.

\subsection{Stream 2: Inventory Management under Power Demand}

\subsubsection*{Sicilia, Gonz\'alez-De-la-Rosa, Febles-Acosta \& Alcaide-L\'opez-de-Pablo (2014)}

The paper of Sicilia and colleagues anchors the second stream. They
analyse a deterministic single-item production-inventory system in which
demand follows the \emph{power demand pattern}
$D(t) = (r/T)\,(t/T)^{(1/n)-1}/n$, parameterised by the index $n>0$.
Values of $n$ greater than one represent demand that
intensifies as the cycle progresses (e.g.,\ petrol or staple
groceries running short toward the end of a replenishment cycle), while
$n<1$ captures demand that decays (e.g.,\ perishable items losing
appeal as the expiry date approaches). Production rate is taken to be
proportional to demand rate, $P(t)=\alpha D(t)$ with $\alpha>1$, so that
inventory accumulates during the replenishment phase, after which a
shortage-and-backlog phase concludes the cycle.

The authors minimise the cycle-averaged sum of holding, shortage, and
ordering costs over two decision variables: the reorder point $s$ and
the scheduling period $T$. They prove that the first-order optimality
condition reduces to a single non-linear equation in a normalised
variable $x \in (0, 1-1/\alpha)$, demonstrate uniqueness of the root via
the Intermediate Value Theorem, and propose a four-step algorithm to
recover the economic lot size $Q^{*}$ and the minimum cost $C^{*}$.
Two limiting cases are recovered as sanity checks: (i)~$n=1$ collapses
the model to the classical economic production quantity (EPQ) with
backorders, and (ii)~$\alpha\to\infty$ recovers the EOQ with power
demand and instantaneous replenishment of Sicilia et al.\ (2012). The
authors deliberately omit any environmental considerations, identifying
deterioration and lost-sales settings as natural extensions, but
\emph{not} carbon emissions.

% ============================================================
\section{Synthesis and Cross-Cutting Observations}
% ============================================================

Three cross-cutting observations emerge from a comparative reading of
the five papers.

\textbf{First, demand is consistently simplified in carbon-aware
models.} Hua et al.\ (2011), Benjaafar et al.\ (2013), and Hasan et al.\
(2021) all retain a \emph{constant-rate} demand assumption (or, at most,
a static linear function of decision variables in Hasan et al.). None of
these models allows demand to vary \emph{within} an inventory cycle,
even though much of the empirical inventory-management literature
documents pronounced intra-cycle non-stationarity for many product
categories. This simplification, while analytically convenient, may
yield biased order-quantity recommendations whenever the underlying
demand exhibits a strong power-form profile.

\textbf{Second, the power-demand literature has remained
environmentally agnostic.} Sicilia et al.\ (2014), and the antecedent
work they cite (Naddor, 1966; Datta and Pal, 1988; Lee and Wu, 2002;
Sicilia et al., 2012; Mishra and Singh, 2013), focus exclusively on
cost minimisation. Neither carbon-emission constraints nor green-
technology investment have been embedded into the power-demand framework.

\textbf{Third, regulatory variety is unevenly modelled across
streams.} The carbon-aware stream has progressed from a single mechanism
(cap-and-trade in Hua et al.) to multi-mechanism comparisons (Benjaafar
et al.\ and Hasan et al., who jointly cover strict caps, taxes,
cap-and-trade, and cap-and-offset). The power-demand stream offers no
analogue. Hence, even if a power-demand inventory model were augmented
with emissions, it would not be obvious \emph{which} regulatory regime
should accompany it.

% ============================================================
\section{Research Gap}
% ============================================================

Synthesising the above, we identify the following research gap, which
motivates further investigation.

\begin{quote}
\textit{No existing study formulates an inventory model that
simultaneously accommodates a power-form, time-varying demand pattern
and one or more carbon-emission regulatory mechanisms (carbon tax,
cap-and-trade, or strict cap), while also allowing for green-technology
investment, backlogged shortages, and a finite production rate.}
\end{quote}

\noindent
This gap can be unpacked into the following more specific deficiencies.

\begin{enumerate}[label=\textbf{G\arabic*.},leftmargin=2.4em,itemsep=0.4em]

\item \textbf{Demand realism in carbon-aware inventory models.}
The three carbon-aware studies reviewed assume constant or
linear-in-investment demand. Extending the power-demand pattern of
Sicilia et al.\ (2014) into the EOQ formulations of Hua et al.\ (2011)
and Hasan et al.\ (2021) would (i) yield richer, more realistic
order-quantity policies and (ii) test whether the central
managerial insights of the carbon-aware literature, e.g.,\ the dominance
of operational adjustment over technology investment in
Benjaafar et al.\ (2013), survive under non-stationary demand.

\item \textbf{Carbon footprint of shortage-and-backlog regimes.}
Sicilia et al.\ (2014) explicitly permit fully backlogged shortages, but
no carbon-aware paper considers the emissions implications of
backordering, e.g.,\ expedited shipping or rush-mode transportation
required to clear a backlog. A model that prices the
\emph{emissions cost of backlogged shortages}, in addition to the usual
monetary shortage cost, would constitute a novel contribution.

\item \textbf{Joint optimisation of order quantity, scheduling period, and green-technology investment.}
The carbon-aware models reviewed optimise either order quantity (Hua
et al.) or order quantity together with a scalar investment level
(Hasan et al.); they do not jointly optimise an additional scheduling
period, as is natural under a time-varying demand. Sicilia et al.\
(2014), by contrast, optimise the scheduling period but not the
investment level. A unified model that treats $Q$, $T$, and $G$ as
co-decision variables remains absent.

\item \textbf{Regulatory comparison under non-stationary demand.}
Hasan et al.\ (2021) and Benjaafar et al.\ (2013) compare regulatory
mechanisms only under static demand. Whether the qualitative ranking of
mechanisms, e.g.,\ Hasan et al.'s observation that the carbon tax rate
matters more than the cap level, persists when demand follows a power
pattern is an open empirical question with policy implications.

\item \textbf{Demand-investment coupling under power demand.}
Hasan et al.\ (2021) introduce demand as an increasing function of
green-technology investment and promotion level. No corresponding
formulation exists in the power-demand literature; in particular, the
\emph{shape parameter} $n$ of the power pattern itself may plausibly
depend on green-marketing effort, an avenue that has been entirely
overlooked.

\item \textbf{Deteriorating items at the intersection.} Sicilia et al.\
(2014) explicitly flag deterioration as a future direction for the
power-demand model, while none of the carbon-aware papers reviewed
incorporates deterioration. A model that combines deterioration, power
demand, and a carbon-pricing mechanism would address two open
directions simultaneously.

\end{enumerate}

\noindent\textbf{Practical Implications.}\;
Closing this gap would benefit several stakeholders. Operations managers
of products with strong power-form demand (e.g.\ fuel retailers,
staple-food distributors, fresh-produce wholesalers) currently lack
analytical tools that calibrate replenishment and shortage-management
decisions to a carbon-priced environment. Policy makers, who already
ask whether a carbon tax or a cap-and-trade scheme is preferable, would
gain insight into how the answer is modulated by the demand profile of
regulated industries. Finally, sustainability-oriented firms
contemplating green-technology investment would acquire a sharper view of
the conditions under which such investment yields net profit improvement.

% ============================================================
\section{Conclusion}
% ============================================================

The five papers reviewed here collectively map two well-developed but
largely separate research streams. The carbon-aware EOQ stream,
represented by Hua, Cheng \& Wang (2011), Benjaafar, Li \& Daskin
(2013), and Hasan, Roy, Daryanto \& Wee (2021), has progressively
broadened the set of regulatory mechanisms, decision variables (now
including green-technology investment), and behavioural ingredients
(promotion-sensitive demand). The non-stationary-demand stream,
exemplified by Sicilia, Gonz\'alez-De-la-Rosa, Febles-Acosta \&
Alcaide-L\'opez-de-Pablo (2014), has refined the algorithmic toolkit
for power-demand systems. Their integration, an inventory model that
simultaneously features a power-form demand pattern, multiple carbon-
emission regulatory mechanisms, optional green-technology investment,
and backlogged shortages, constitutes a clear and as-yet-unaddressed
research opportunity at the frontier of sustainable operations.

% ============================================================
\section*{References}
\addcontentsline{toc}{section}{References}
% ============================================================

\begin{enumerate}[leftmargin=2em,itemsep=0.35em,label={[\arabic*]}]
\item Benjaafar, S., Li, Y., \& Daskin, M.\ (2013). Carbon Footprint and the Management of Supply Chains: Insights From Simple Models. \emph{IEEE Transactions on Automation Science and Engineering}, 10(1), 99--116.

\item Hasan, M.\,R., Roy, T.\,C., Daryanto, Y., \& Wee, H.-M.\ (2021). Optimizing inventory level and technology investment under a carbon tax, cap-and-trade and strict carbon limit regulations. \emph{Sustainable Production and Consumption}, 25, 604--621.

\item Hua, G., Cheng, T.\,C.\,E., \& Wang, S.\ (2011). Managing carbon footprints in inventory management. \emph{International Journal of Production Economics}, 132(2), 178--185.

\item Hua, G., Cheng, T.\,C.\,E., \& Wang, S.\ (2010). Managing carbon footprints in inventory management. SSRN Working Paper No.\ 1628953. (Pre-print of Hua et al., 2011.)

\item Sicilia, J., Gonz\'alez-De-la-Rosa, M., Febles-Acosta, J., \& Alcaide-L\'opez-de-Pablo, D.\ (2014). Optimal policy for an inventory system with power demand, backlogged shortages and production rate proportional to demand rate. \emph{International Journal of Production Economics}, 155, 163--171.
\end{enumerate}

\end{document}
